% ============================================================
% LECCIÓN: FLUIDOS
% Curso de Oriente - CCH Oriente
% Enero 2026
% ============================================================
% INSTRUCCIONES: Pegar este contenido en tu plantilla de Overleaf
% después de la definición de entornos y antes de \end{document}
% ============================================================

\renewcommand{\tema}{Fluidos}

% ============================================================
\section{Fluidos en Reposo}
% ============================================================

\begin{seccion}
Un \textbf{fluido} es toda sustancia capaz de fluir y adoptar la forma del recipiente que lo contiene. Comprende líquidos y gases. A diferencia de los sólidos, los fluidos no resisten la deformación por cortante: basta aplicar la mínima fuerza tangencial para que la forma cambie. En este bloque estudiaremos el comportamiento de los fluidos cuando están en equilibrio estático.
\end{seccion}

% ============================================================
\subsection{Presión}
% ============================================================

\begin{seccion}
La presión cuantifica cómo se distribuye una fuerza sobre una superficie. Dos fuerzas iguales pueden producir efectos muy distintos según el área sobre la que actúan: pisar con el tacón de un zapato duele más que pisar con la planta, aunque el peso sea el mismo.

\begin{definicion}
\textbf{Presión ($P$):} cociente entre la fuerza perpendicular aplicada sobre una superficie y el área de esa superficie.
\end{definicion}

\begin{ecuacion}
\textbf{Definición de presión:}
$$P = \frac{F}{A}$$

donde:
\begin{itemize}
    \item $P$ = presión (Pa $= \dfrac{\text{N}}{\text{m}^2}$)
    \item $F$ = fuerza perpendicular a la superficie (N)
    \item $A$ = área de la superficie (m$^2$)
\end{itemize}
\end{ecuacion}

\begin{nota}
\textbf{Unidades de presión más comunes:}
\begin{itemize}
    \item 1 Pa $= 1$ N/m$^2$ (unidad SI)
    \item 1 atm $= 101{,}325$ Pa $\approx 101$ kPa
    \item 1 atm $= 760$ mmHg $= 760$ torr
    \item 1 bar $= 100{,}000$ Pa
\end{itemize}
La presión es una magnitud \textbf{escalar}: no tiene dirección.
\end{nota}
\end{seccion}

% ============================================================
\subsection{Presión Atmosférica}
% ============================================================

\begin{seccion}
La atmósfera terrestre es una capa de gases que rodea el planeta. Su peso ejerce presión sobre todo lo que se encuentra en la superficie. Esta presión, llamada \textbf{presión atmosférica}, equivale al peso de una columna de aire desde el nivel del mar hasta el límite de la atmósfera.

\begin{definicion}
\textbf{Presión atmosférica estándar ($P_0$):} presión ejercida por la atmósfera al nivel del mar.
$$P_0 = 1\text{ atm} = 101{,}325\text{ Pa} \approx 101\text{ kPa}$$
\end{definicion}

El físico Evangelista Torricelli demostró en 1644 que la presión atmosférica es capaz de sostener una columna de mercurio de exactamente 760 mm. Este principio es la base del \textbf{barómetro}.

\begin{nota}
La presión atmosférica \textbf{disminuye con la altitud} porque hay menos columna de aire encima. En la Ciudad de México (2\,240 m sobre el nivel del mar), la presión es de aproximadamente 79 kPa, equivalente a $\approx 0.78$ atm.
\end{nota}
\end{seccion}

% ============================================================
\subsection{Presión Hidrostática}
% ============================================================

\begin{seccion}
Al sumergirse en un fluido, la presión aumenta con la profundidad debido al peso del fluido que hay encima. Esta presión debida a la columna de fluido se llama \textbf{presión hidrostática}.

\begin{ecuacion}
\textbf{Presión hidrostática:}
$$P = P_0 + \rho g h$$

donde:
\begin{itemize}
    \item $P$ = presión total a la profundidad $h$ (Pa)
    \item $P_0$ = presión en la superficie del fluido (Pa)
    \item $\rho$ = densidad del fluido (kg/m$^3$)
    \item $g$ = aceleración gravitacional ($\approx 10$ m/s$^2$)
    \item $h$ = profundidad medida desde la superficie libre (m)
\end{itemize}
\end{ecuacion}

\textbf{Consecuencias importantes:}
\begin{itemize}
    \item La presión hidrostática \textbf{no depende} de la forma del recipiente ni de la cantidad total de fluido, solo de la profundidad y la densidad.
    \item A la misma profundidad, la presión es igual en todos los puntos de un fluido en reposo (\textbf{principio de isobaricidad horizontal}).
    \item La diferencia entre la presión total y la presión atmosférica se llama \textbf{presión manométrica}: $P_{\text{man}} = \rho g h$.
\end{itemize}
\end{seccion}

\subsubsection{Diagrama: Presión en función de la profundidad}

\begin{center}
\begin{tikzpicture}[scale=1.0]

    % Recipiente
    \draw[thick, azuloscuro] (0,0) rectangle (4,5);
    % Agua
    \fill[blue!15] (0.05,0.05) rectangle (3.95,4.95);

    % Superficie libre
    \draw[dashed, azulmedio, thick] (0,5) -- (5.8,5);
    \node[right, azulmedio] at (5.8,5) {Superficie libre ($P_0$)};

    % Línea h1
    \draw[dashed, gray] (0,3.5) -- (5.8,3.5);
    \draw[thick, azuloscuro, -{Stealth}] (4.2,3.5) -- (5.1,3.5);
    \node[right] at (5.1,3.5) {presión menor};

    % Línea h2
    \draw[dashed, gray] (0,1.5) -- (5.8,1.5);
    \draw[thick, azuloscuro, -{Stealth}] (4.2,1.5) -- (5.1,1.5);
    \node[right] at (5.1,1.5) {presión mayor};

    % Cotas de profundidad
    \draw[{Stealth}-{Stealth}, gray] (-0.6,5) -- (-0.6,3.5);
    \node[left, gray] at (-0.6,4.25) {$h_1$};
    \draw[{Stealth}-{Stealth}, gray] (-0.6,5) -- (-0.6,1.5);
    \node[left, gray] at (-0.6,3.25) {$h_2$};

    % Etiqueta fórmula y fluido
    \node[azuloscuro] at (2,2.5) {\large fluido ($\rho$)};
    \node[right] at (4.2,0.5) {\small $P = P_0 + \rho g h$};

\end{tikzpicture}
\end{center}

% ============================================================
\subsection{Principio de Pascal}
% ============================================================

\begin{seccion}
Blaise Pascal enunció en el siglo XVII uno de los principios más útiles de la hidrostática, que es el fundamento de toda la hidráulica moderna.

\begin{definicion}
\textbf{Principio de Pascal:} toda variación de presión aplicada a un fluido confinado e incompresible se transmite íntegramente a todos los puntos del fluido y a las paredes del recipiente.
\end{definicion}

La aplicación más directa de este principio es la \textbf{prensa hidráulica}. Si se aplica una fuerza $F_1$ sobre un pistón de área $A_1$, la presión generada es $P = F_1/A_1$. Como esa presión se transmite sin pérdida al segundo pistón de área $A_2$, la fuerza resultante es $F_2 = P \cdot A_2$.

\begin{ecuacion}
\textbf{Prensa hidráulica:}
$$\frac{F_1}{A_1} = \frac{F_2}{A_2} \qquad \Longrightarrow \qquad F_2 = F_1 \cdot \frac{A_2}{A_1}$$

Si $A_2 > A_1$, entonces $F_2 > F_1$: la prensa \textbf{multiplica} la fuerza.
\end{ecuacion}

\textbf{Aplicaciones del principio de Pascal:}
\begin{itemize}
    \item Frenos hidráulicos de automóviles
    \item Gato hidráulico para cambiar llantas
    \item Prensas industriales y elevadores de taller
    \item Dirección hidráulica asistida
\end{itemize}
\end{seccion}

\subsubsection{Diagrama: Prensa hidráulica}

\begin{center}
\begin{tikzpicture}[scale=1.0, line join=round]

    % --- Fluido ---
    % Canal de conexión inferior
    \fill[blue!18] (0.5,0) rectangle (6.0,0.8);
    % Columna izquierda (hasta el pistón en y=2.0)
    \fill[blue!18] (0.5,0.8) rectangle (2.0,2.0);
    % Columna derecha (hasta el pistón en y=2.8)
    \fill[blue!18] (3.5,0.8) rectangle (6.0,2.8);

    % --- Paredes del cilindro izquierdo (angosto) ---
    \draw[thick, azuloscuro] (0.5,0) -- (0.5,4.5);
    \draw[thick, azuloscuro] (2.0,0) -- (2.0,4.5);

    % --- Paredes del cilindro derecho (ancho) ---
    \draw[thick, azuloscuro] (3.5,0) -- (3.5,4.5);
    \draw[thick, azuloscuro] (6.0,0) -- (6.0,4.5);

    % --- Canal inferior ---
    \draw[thick, azuloscuro] (0.5,0) -- (6.0,0);       % base
    \draw[thick, azuloscuro] (2.0,0.8) -- (3.5,0.8);   % techo del canal

    % --- Pistón izquierdo ---
    \fill[gray!65] (0.5,2.0) rectangle (2.0,2.35);
    \draw[thick] (0.5,2.0) rectangle (2.0,2.35);

    % --- Pistón derecho ---
    \fill[gray!65] (3.5,2.8) rectangle (6.0,3.15);
    \draw[thick] (3.5,2.8) rectangle (6.0,3.15);

    % --- Fuerza F1 (hacia abajo) ---
    \draw[very thick, -{Stealth}, verdecorrecto] (1.25,3.7) -- (1.25,2.35);
    \node[above, verdecorrecto] at (1.25,3.8) {$F_1$};

    % --- Fuerza F2 (hacia arriba) ---
    \draw[very thick, -{Stealth}, azuloscuro] (4.75,3.15) -- (4.75,4.5);
    \node[above, azuloscuro] at (4.75,4.6) {$F_2$};

    % --- Cotas de área ---
    \draw[{Stealth}-{Stealth}, gray, thin] (0.5,-0.6) -- (2.0,-0.6);
    \node[below, gray] at (1.25,-0.6) {\small $A_1$ (pequeña)};

    \draw[{Stealth}-{Stealth}, gray, thin] (3.5,-0.6) -- (6.0,-0.6);
    \node[below, gray] at (4.75,-0.6) {\small $A_2$ (grande)};

    % --- Etiqueta de fluido ---
    \node[azuloscuro!80] at (3.05,0.4) {\small fluido};

\end{tikzpicture}
\end{center}

% ============================================================
\subsection{Principio de Arquímedes}
% ============================================================

\begin{seccion}
Cuando un objeto se sumerge total o parcialmente en un fluido, el fluido ejerce sobre él una fuerza hacia arriba llamada \textbf{empuje} o \textbf{fuerza boyante}. Arquímedes de Siracusa enunció el principio que permite calcularla.

\begin{definicion}
\textbf{Principio de Arquímedes:} todo cuerpo sumergido en un fluido experimenta una fuerza de empuje dirigida hacia arriba, igual al peso del fluido desplazado por el cuerpo.
\end{definicion}

\begin{ecuacion}
\textbf{Fuerza de empuje:}
$$F_e = \rho_{\text{fluido}} \cdot g \cdot V_{\text{sumergido}}$$

donde:
\begin{itemize}
    \item $F_e$ = fuerza de empuje (N), dirigida hacia arriba
    \item $\rho_{\text{fluido}}$ = densidad del fluido (kg/m$^3$)
    \item $g$ = aceleración gravitacional ($\approx 10$ m/s$^2$)
    \item $V_{\text{sumergido}}$ = volumen del objeto que está dentro del fluido (m$^3$)
\end{itemize}
\end{ecuacion}

\textbf{Condiciones de flotación:}
\begin{itemize}
    \item Si $\rho_{\text{obj}} < \rho_{\text{fluido}}$: el objeto \textbf{flota}.
    \item Si $\rho_{\text{obj}} = \rho_{\text{fluido}}$: el objeto queda en \textbf{equilibrio} a cualquier profundidad.
    \item Si $\rho_{\text{obj}} > \rho_{\text{fluido}}$: el objeto se \textbf{hunde}.
\end{itemize}

Cuando un objeto de densidad $\rho_{\text{obj}}$ flota sobre un fluido de densidad $\rho_{\text{fluido}}$, la fracción de su volumen bajo la superficie es:

$$\frac{V_{\text{sumergido}}}{V_{\text{total}}} = \frac{\rho_{\text{obj}}}{\rho_{\text{fluido}}}$$
\end{seccion}

\subsubsection{Diagrama: Fuerza de empuje sobre un objeto sumergido}

\begin{center}
\begin{tikzpicture}[scale=1.0]

    % Recipiente con fluido
    \fill[blue!12] (0,0) rectangle (6,5);
    \draw[thick, azuloscuro] (0,0) rectangle (6,5);
    \draw[dashed, azulmedio] (0,5) -- (6.8,5);
    \node[right, azulmedio] at (6.8,5) {Superficie};

    % Objeto sumergido (cubo)
    \fill[gray!50] (2.2,1.8) rectangle (3.8,3.2);
    \draw[thick] (2.2,1.8) rectangle (3.8,3.2);
    \node at (3.0,2.5) {\small objeto};

    % Peso (hacia abajo)
    \draw[very thick, -{Stealth}, red] (3.0,1.8) -- (3.0,0.4);
    \node[right, red] at (3.05,1.1) {$W = mg$};

    % Empuje (hacia arriba)
    \draw[very thick, -{Stealth}, verdecorrecto] (3.0,3.2) -- (3.0,4.7);
    \node[right, verdecorrecto] at (3.05,3.9) {$F_e = \rho_f g V$};

    % Densidad del fluido
    \node[azuloscuro, opacity=0.6] at (1.0,2.5) {\Large $\rho_f$};
    \node[azuloscuro, opacity=0.6] at (5.0,2.5) {\Large $\rho_f$};

\end{tikzpicture}
\end{center}

% ============================================================
\subsection{Tensión Superficial y Capilaridad}
% ============================================================

\begin{seccion}
En la interfaz entre un líquido y el aire, las moléculas del líquido solo interactúan con otras moléculas del líquido por debajo y a los lados; no hay moléculas de líquido arriba. Esto genera una fuerza neta hacia el interior que hace que la superficie se comporte como una membrana elástica bajo tensión.

\begin{definicion}
\textbf{Tensión superficial ($\gamma$):} fuerza por unidad de longitud que actúa sobre la superficie de un líquido y tiende a minimizar su área superficial. Se mide en N/m.
\end{definicion}

\textbf{Manifestaciones cotidianas:}
\begin{itemize}
    \item Las gotas de agua adoptan forma esférica.
    \item Los insectos como el zapatero caminan sobre la superficie del agua sin hundirse.
    \item Una aguja de acero puede flotar horizontalmente sobre el agua.
    \item El agua de lluvia forma esferas sobre tela impermeable.
\end{itemize}

\vspace{0.3cm}

\textbf{Capilaridad:}

Cuando un tubo de diámetro muy pequeño se pone en contacto con un líquido, el líquido puede ascender o descender por el tubo. Este fenómeno se llama \textbf{capilaridad} y resulta de la competencia entre la \textbf{cohesión} (atracción entre moléculas del mismo líquido) y la \textbf{adhesión} (atracción entre el líquido y el sólido).

\begin{itemize}
    \item Si la adhesión supera la cohesión (agua en vidrio): el líquido \textbf{asciende} y el menisco es cóncavo.
    \item Si la cohesión supera la adhesión (mercurio en vidrio): el líquido \textbf{desciende} y el menisco es convexo.
    \item A menor diámetro del tubo, mayor es el efecto capilar.
\end{itemize}

\textbf{Importancia biológica y tecnológica:}
\begin{itemize}
    \item Ascenso de savia en los árboles hasta decenas de metros de altura.
    \item Absorción de agua por las raíces a través del suelo.
    \item Funcionamiento de telas y materiales absorbentes.
\end{itemize}
\end{seccion}

\subsubsection{Diagrama: Capilaridad — tubos con menisco cóncavo y convexo}

\begin{center}
\begin{tikzpicture}[scale=1.0]

    % Recipiente
    \fill[blue!12] (-3.0,-1.2) rectangle (4.5,0);
    \draw[thick, azuloscuro] (-3.0,-1.2) rectangle (4.5,0.05);
    \draw[dashed, gray] (-3.0,0) -- (4.8,0);
    \node[right, gray] at (4.8,0) {\small nivel del líquido};

    % --- Tubo capilar izquierdo: agua en vidrio (asciende) ---
    \draw[thick, azuloscuro] (-0.8,0) -- (-0.8,3.8);
    \draw[thick, azuloscuro] (0.8,0)  -- (0.8,3.8);
    % Fluido dentro del tubo
    \fill[blue!30] (-0.79,0) rectangle (0.79,2.6);
    % Menisco cóncavo
    \draw[thick, blue!70] (-0.79,2.6) .. controls (0,2.95) .. (0.79,2.6);
    % Cota h
    \draw[{Stealth}-{Stealth}, gray, thin] (-1.2,0) -- (-1.2,2.6);
    \node[left, gray] at (-1.2,1.3) {$h$};
    % Etiquetas
    \node[above] at (0,3.0) {\small menisco cóncavo};
    \node[below] at (0,-1.5) {\small agua en vidrio};
    \node[right] at (0.85,1.3) {\small \textcolor{azuloscuro}{asciende}};

    % --- Tubo capilar derecho: mercurio en vidrio (desciende) ---
    \draw[thick, azuloscuro] (2.2,0) -- (2.2,3.8);
    \draw[thick, azuloscuro] (3.8,0) -- (3.8,3.8);
    % Fluido dentro del tubo (nivel más bajo que el recipiente)
    \fill[gray!50] (2.21,0) rectangle (3.79,-0.6);
    % Menisco convexo
    \draw[thick, gray!70] (2.21,0) .. controls (3.0,-0.35) .. (3.79,0);
    % Etiquetas
    \node[above] at (3.0,0.15) {\small menisco convexo};
    \node[below] at (3.0,-1.5) {\small mercurio en vidrio};
    \node[right] at (3.85,-0.3) {\small \textcolor{gray}{desciende}};

\end{tikzpicture}
\end{center}

% ============================================================
\section{Fluidos en Movimiento}
% ============================================================

\begin{seccion}
Cuando un fluido se desplaza de un punto a otro, se dice que está en régimen de \textbf{flujo}. Para simplificar el análisis, trabajaremos con fluidos \textbf{ideales}: incompresibles (densidad constante), no viscosos (sin fricción interna) y en flujo laminar (capas ordenadas sin mezcla). Estas aproximaciones son válidas para muchas situaciones prácticas con líquidos.
\end{seccion}

% ============================================================
\subsection{Gasto y Flujo}
% ============================================================

\begin{seccion}

\begin{definicion}
\textbf{Gasto volumétrico ($Q$):} volumen de fluido que pasa por una sección transversal por unidad de tiempo.
\end{definicion}

\begin{ecuacion}
\textbf{Gasto volumétrico:}
$$Q = A \cdot v$$

donde:
\begin{itemize}
    \item $Q$ = gasto volumétrico (m$^3$/s)
    \item $A$ = área de la sección transversal (m$^2$)
    \item $v$ = velocidad del fluido (m/s)
\end{itemize}
\end{ecuacion}

\textbf{Tipos de flujo:}
\begin{itemize}
    \item \textbf{Flujo laminar:} el fluido se mueve en capas paralelas bien definidas, sin mezcla entre ellas. Se presenta a velocidades bajas.
    \item \textbf{Flujo turbulento:} el movimiento es caótico, con remolinos y mezcla entre capas. Se presenta a velocidades altas.
\end{itemize}

\end{seccion}

% ============================================================
\subsection{Viscosidad}
% ============================================================

\begin{seccion}

\begin{definicion}
\textbf{Viscosidad:} propiedad de un fluido que mide su resistencia a fluir, es decir, la fricción interna entre sus capas. A mayor viscosidad, el fluido fluye más lentamente ante la misma fuerza aplicada.
\end{definicion}

\begin{nota}
\textbf{Ejemplos comparativos de viscosidad (de menor a mayor):}

\hspace{1cm} Aire $<$ agua $<$ aceite vegetal $<$ miel $<$ brea

La viscosidad del agua \textbf{disminuye} al aumentar la temperatura; la de los gases \textbf{aumenta}. En fluidos ideales, la viscosidad se considera nula.
\end{nota}

\end{seccion}

% ============================================================
\subsection{Ecuación de Continuidad}
% ============================================================

\begin{seccion}
Si un fluido incompresible fluye por un tubo sin fugas ni inyecciones, la cantidad de fluido que entra por un extremo es exactamente la misma que sale por el otro. Esta es una consecuencia de la \textbf{conservación de la masa}.

\begin{ecuacion}
\textbf{Ecuación de continuidad:}
$$A_1 v_1 = A_2 v_2 \qquad \Longrightarrow \qquad Q = \text{constante}$$

donde:
\begin{itemize}
    \item $A_1$, $A_2$ = áreas de la sección transversal en los puntos 1 y 2 (m$^2$)
    \item $v_1$, $v_2$ = velocidades del fluido en los puntos 1 y 2 (m/s)
\end{itemize}
\end{ecuacion}

\textbf{Consecuencia fundamental:} donde el tubo es más angosto ($A$ disminuye), el fluido se mueve más rápido ($v$ aumenta), y viceversa. Cuando uno aprieta la boquilla de una manguera, el chorro sale más veloz.
\end{seccion}

\subsubsection{Diagrama: Ecuación de continuidad — tubo con sección variable}

\begin{center}
\begin{tikzpicture}[scale=1.0]

    % Relleno del tubo (fluido)
    \fill[blue!15]
        (0,0.8) -- (3,0.8) -- (4.5,0.2) -- (8,0.2) --
        (8,-0.2) -- (4.5,-0.2) -- (3,-0.8) -- (0,-0.8) -- cycle;

    % Paredes del tubo
    \draw[thick, azuloscuro] (0,0.8)  -- (3,0.8)  -- (4.5,0.2)  -- (8,0.2);
    \draw[thick, azuloscuro] (0,-0.8) -- (3,-0.8) -- (4.5,-0.2) -- (8,-0.2);

    % Sección 1 (sección ancha)
    \draw[dashed, gray] (1.2,0.8) -- (1.2,-0.8);
    \draw[{Stealth}-{Stealth}, gray] (1.2,0.9) -- (1.2,-0.9);
    \node[above, gray] at (1.2,1.0) {$A_1$};

    % Sección 2 (sección angosta)
    \draw[dashed, gray] (6.5,0.2) -- (6.5,-0.2);
    \draw[{Stealth}-{Stealth}, gray] (6.5,0.3) -- (6.5,-0.3);
    \node[above, gray] at (6.5,0.4) {$A_2$};

    % Vector velocidad 1
    \draw[very thick, -{Stealth}, verdecorrecto] (0.1,0) -- (1.0,0);
    \node[above, verdecorrecto] at (0.55,0.05) {\small $v_1$};

    % Vector velocidad 2 (más largo, mayor velocidad)
    \draw[very thick, -{Stealth}, azuloscuro] (5.8,0) -- (7.6,0);
    \node[above, azuloscuro] at (6.7,0.05) {\small $v_2 > v_1$};

    % Nota A1 > A2
    \node[below] at (4.0,-1.1) {$A_1 > A_2 \;\Rightarrow\; v_2 > v_1$};

\end{tikzpicture}
\end{center}

% ============================================================
\subsection{Ecuación de Bernoulli}
% ============================================================

\begin{seccion}
Daniel Bernoulli demostró en 1738 que en un fluido ideal en movimiento, la suma de la presión, la energía cinética por unidad de volumen y la energía potencial gravitacional por unidad de volumen se conserva a lo largo de una línea de flujo. Es la \textbf{conservación de la energía} aplicada a fluidos.

\begin{ecuacion}
\textbf{Ecuación de Bernoulli:}
$$P_1 + \frac{1}{2}\rho v_1^2 + \rho g h_1 = P_2 + \frac{1}{2}\rho v_2^2 + \rho g h_2 = \text{cte}$$

donde:
\begin{itemize}
    \item $P$ = presión en el punto considerado (Pa)
    \item $\rho$ = densidad del fluido (kg/m$^3$)
    \item $v$ = velocidad del fluido (m/s)
    \item $g$ = aceleración gravitacional ($\approx 10$ m/s$^2$)
    \item $h$ = altura medida desde un nivel de referencia (m)
\end{itemize}
\end{ecuacion}

\textbf{Casos especiales importantes:}

\begin{itemize}
    \item \textbf{Tubo horizontal} ($h_1 = h_2$): se elimina el término gravitacional.
    $$P_1 + \frac{1}{2}\rho v_1^2 = P_2 + \frac{1}{2}\rho v_2^2$$
    Donde el fluido va más rápido, la presión es \textbf{menor}.

    \item \textbf{Fluido en reposo} ($v_1 = v_2 = 0$): recuperamos la ecuación de presión hidrostática.
    $$P_1 + \rho g h_1 = P_2 + \rho g h_2$$
\end{itemize}

\textbf{Efecto Venturi:} cuando un fluido pasa por una sección reducida del tubo, su velocidad aumenta (por la ecuación de continuidad) y su presión disminuye (por Bernoulli). Este principio se aprovecha en carburadores, atomizadores y medidores de flujo.
\end{seccion}

\subsubsection{Diagrama: Efecto Venturi}

\begin{center}
\begin{tikzpicture}[scale=1.0]

    % Fluido dentro del tubo venturi
    \fill[blue!15]
        (0,0.8) -- (2.5,0.8) -- (3.5,0.3) -- (5.5,0.3) -- (6.5,0.8) -- (9,0.8) --
        (9,-0.8) -- (6.5,-0.8) -- (5.5,-0.3) -- (3.5,-0.3) -- (2.5,-0.8) -- (0,-0.8) -- cycle;

    % Paredes del tubo
    \draw[thick, azuloscuro]
        (0,0.8) -- (2.5,0.8) -- (3.5,0.3) -- (5.5,0.3) -- (6.5,0.8) -- (9,0.8);
    \draw[thick, azuloscuro]
        (0,-0.8) -- (2.5,-0.8) -- (3.5,-0.3) -- (5.5,-0.3) -- (6.5,-0.8) -- (9,-0.8);

    % Manómetro izquierdo (presión alta, columna alta)
    \draw[thick, gray] (1.5,0.8) -- (1.5,3.2);
    \fill[blue!40] (1.3,0.8) rectangle (1.7,2.8);
    \draw[thick, blue!60] (1.3,2.8) -- (1.7,2.8);
    \node[above] at (1.5,3.2) {\small $P_1$ (alta)};

    % Manómetro central (presión baja, columna baja)
    \draw[thick, gray] (4.5,0.3) -- (4.5,3.2);
    \fill[blue!40] (4.3,0.3) rectangle (4.7,1.2);
    \draw[thick, blue!60] (4.3,1.2) -- (4.7,1.2);
    \node[above] at (4.5,3.2) {\small $P_2$ (baja)};

    % Vectores de velocidad
    \draw[very thick, -{Stealth}, verdecorrecto] (0.2,0) -- (1.1,0);
    \node[above, verdecorrecto] at (0.65,0.05) {\small $v_1$};

    \draw[very thick, -{Stealth}, azuloscuro] (3.7,0) -- (5.3,0);
    \node[above, azuloscuro] at (4.5,0.05) {\small $v_2 > v_1$};

    \draw[very thick, -{Stealth}, verdecorrecto] (7.9,0) -- (8.8,0);
    \node[above, verdecorrecto] at (8.35,0.05) {\small $v_1$};

    % Leyenda inferior
    \node at (4.5,-1.3) {\small Donde $v$ aumenta, $P$ disminuye};

\end{tikzpicture}
\end{center}

% ============================================================
\subsection{Teorema de Torricelli}
% ============================================================

\begin{seccion}
El teorema de Torricelli describe la velocidad con que un líquido sale por un orificio en las paredes o la base de un recipiente abierto. Es un caso particular de la ecuación de Bernoulli en el que la velocidad en la superficie es despreciable respecto a la del chorro.

\textbf{Condiciones de aplicación:}
\begin{itemize}
    \item El recipiente es grande, de modo que la superficie baja muy lentamente: $v_{\text{sup}} \approx 0$.
    \item El orificio está abierto a la atmósfera: $P_{\text{orificio}} = P_0$.
    \item La presión sobre la superficie también es $P_0$ (recipiente abierto).
\end{itemize}

\textbf{Derivación:} Aplicando Bernoulli entre la superficie libre (punto 1) y el orificio (punto 2), tomando la altura del orificio como referencia ($h_2 = 0$, $h_1 = h$):

$$P_0 + \frac{1}{2}\rho (0)^2 + \rho g h \;=\; P_0 + \frac{1}{2}\rho v^2 + \rho g (0)$$

$$\rho g h = \frac{1}{2}\rho v^2 \qquad \Longrightarrow \qquad v^2 = 2gh$$

\begin{ecuacion}
\textbf{Teorema de Torricelli:}
$$v = \sqrt{2gh}$$

donde:
\begin{itemize}
    \item $v$ = velocidad del fluido al salir por el orificio (m/s)
    \item $g$ = aceleración gravitacional ($\approx 10$ m/s$^2$)
    \item $h$ = altura de la superficie libre sobre el orificio (m)
\end{itemize}
\end{ecuacion}

\begin{nota}
\textbf{Analogía notable:} la expresión $v = \sqrt{2gh}$ es idéntica a la velocidad que adquiere un objeto en caída libre desde la misma altura $h$. Torricelli observó esta equivalencia antes de que Newton formulara las leyes del movimiento.
\end{nota}
\end{seccion}

\subsubsection{Diagrama: Tanque con orificio lateral — Torricelli}

\begin{center}
\begin{tikzpicture}[scale=1.0]

    % Tanque con fluido
    \fill[blue!15] (0,0) rectangle (4,5);
    % Paredes del tanque (sin el orificio)
    \draw[thick, azuloscuro] (0,0) -- (0,5) -- (4,5) -- (4,2.2);
    \draw[thick, azuloscuro] (4,1.8) -- (4,0) -- (0,0);

    % Chorro de agua saliendo por el orificio
    \fill[blue!35]
        (4.0,1.85) .. controls (5.2,1.85) and (6.2,0.9) .. (6.8,0.4) --
        (6.8,0.2) .. controls (6.2,0.7) and (5.2,1.65) .. (4.0,1.65) -- cycle;

    % Flecha de velocidad
    \draw[very thick, -{Stealth}, azuloscuro] (4.2,1.75) -- (5.6,1.75);
    \node[above, azuloscuro] at (4.9,1.85) {$v = \sqrt{2gh}$};

    % Superficie libre
    \draw[dashed, azulmedio] (0,5) -- (5.0,5);
    \node[right, azulmedio] at (5.0,5) {sup. libre ($v \approx 0$)};

    % Cota h
    \draw[{Stealth}-{Stealth}, gray] (4.7,1.75) -- (4.7,5.0);
    \node[right, gray] at (4.8,3.4) {$h$};

    % Etiquetas de puntos
    \node[blue!70] at (2,5.35) {\small punto 1: $P_0,\; v \approx 0$};
    \node[blue!70] at (4.0,1.35) {\small punto 2: $P_0$};

\end{tikzpicture}
\end{center}

% ============================================================
\section*{Formulario}
% ============================================================

\begin{nota}

\textbf{Presión:}
$$P = \frac{F}{A} \qquad [P] = \text{Pa} = \frac{\text{N}}{\text{m}^2}$$

\vspace{0.3cm}

\textbf{Presión hidrostática:}
$$P = P_0 + \rho g h$$

\vspace{0.3cm}

\textbf{Principio de Pascal — Prensa hidráulica:}
$$\frac{F_1}{A_1} = \frac{F_2}{A_2}$$

\vspace{0.3cm}

\textbf{Principio de Arquímedes — Fuerza de empuje:}
$$F_e = \rho_{\text{fluido}} \cdot g \cdot V_{\text{sumergido}}$$

\vspace{0.3cm}

\textbf{Fracción sumergida en equilibrio de flotación:}
$$\frac{V_{\text{sumergido}}}{V_{\text{total}}} = \frac{\rho_{\text{obj}}}{\rho_{\text{fluido}}}$$

\vspace{0.3cm}

\textbf{Gasto volumétrico:}
$$Q = A \cdot v \qquad [Q] = \text{m}^3/\text{s}$$

\vspace{0.3cm}

\textbf{Ecuación de continuidad:}
$$A_1 v_1 = A_2 v_2$$

\vspace{0.3cm}

\textbf{Ecuación de Bernoulli:}
$$P_1 + \frac{1}{2}\rho v_1^2 + \rho g h_1 = P_2 + \frac{1}{2}\rho v_2^2 + \rho g h_2$$

\vspace{0.3cm}

\textbf{Teorema de Torricelli:}
$$v = \sqrt{2gh}$$

\vspace{0.3cm}

\textbf{Datos de referencia:}
\begin{itemize}
    \item $g \approx 10$ m/s$^2$
    \item $\rho_{\text{agua}} = 1{,}000$ kg/m$^3$
    \item $\rho_{\text{mercurio}} = 13{,}600$ kg/m$^3$
    \item $\rho_{\text{aire}} \approx 1.29$ kg/m$^3$
    \item $P_0 = 1$ atm $= 101{,}325$ Pa $\approx 101$ kPa $= 760$ mmHg
\end{itemize}

\end{nota}
